\chapter{总结与展望}

\section{工作总结}

本文围绕“基于云原生与AI辅助的一体化在线编程评测平台设计与实现”这一课题，完成了需求分析、总体设计、详细实现与测试分析，形成了面向高校程序设计教学场景的在线评测平台原型。论文的主线目标是把客观判题、容器化执行和 AI 辅助反馈整合到同一平台中。

在实现层面，平台已完成题库管理、题目展示、模板加载、在线提交、Docker 隔离编译运行、提交记录展示和 AI 分析等核心功能。第六章的测试结果表明，系统能够正确返回 \texttt{AC}、\texttt{CE}、\texttt{WA}、\texttt{RE}、\texttt{TLE} 等常见状态，前端工作区、多语言切换、主题切换以及编辑阶段 AI 工具也可以正常使用。项目同时通过了 \texttt{bun run lint} 与 \texttt{bun run build} 校验。

除主线功能外，论文还对异步判题、实时状态回传、课程与作业组织、课程级运行环境配置等扩展方向进行了设计论证。这些内容虽然尚未全部进入当前主线实现，但已经为后续继续迭代保留了结构基础。

\section{不足与展望}

当前系统仍有三个直接限制。第一，主线实现主要围绕 C 和 C++，语言覆盖范围有限；第二，判题流程仍以一次提交对应一次容器执行为主，尚未正式切换到独立 worker 与任务队列结合的异步模式；第三，课程、作业和课程级运行环境配置仍停留在设计和预留阶段。

AI 辅助部分也还有改进空间。当前系统已经能够提供复杂度分析、代码优化建议和问答支持，但在反馈稳定性、教师侧审阅、边界控制和更细粒度学习提示方面还不够完善。后续若继续迭代，可优先补充反馈分级、异常内容过滤以及教师侧摘要能力。

若按更高研究标准要求，本文尚未展开并发提交时延、基础设施成本、环境影响、视觉反馈和学生自编测试评价等实验。这些内容在相关研究中具有独立价值[16-18]，也是当前论文与硕博层面工作之间的主要差距。

从后续实现优先级看，应先补齐课程与作业组织能力、课程级运行环境配置以及异步判题主线；排行榜等展示型功能可以后置；竞赛模块则更适合作为后续独立方向，不宜继续挤占当前毕设主线。
